\section{Metodología}
\label{sec:metodologia}

El desarrollo del proyecto adopta una metodología basada en la 
adaptación del Proceso Unificado Racional (Rational Unified Process, 
RUP) al contexto del Proyecto Integrador, complementada con 
prácticas iterativas e incrementales propias del desarrollo ágil 
moderno. La elección responde a tres características del proyecto: 
la duración semestral con entregables formales por fase, que se 
alinea con el ciclo de vida definido por hitos del RUP; la naturaleza 
iterativa del desarrollo, que requiere ciclos cortos de construcción 
y validación dentro de cada fase; y la necesidad de gestión explícita 
de riesgos técnicos, particularmente en componentes sensibles como la 
captura de tráfico a bajo nivel y el pipeline de análisis distribuido 
\cite{kruchten2003rup}.

\subsection{Adaptación de RUP al Proyecto Integrador}

El RUP organiza el ciclo de vida del software en cuatro fases 
secuenciales que separan el tipo de trabajo predominante en cada 
momento: Inicio, Elaboración, Construcción y Transición 
\cite{jacobson1999up}. Cada fase se desarrolla de forma iterativa 
mediante ciclos numerados 1 a $k$, donde el número de iteraciones 
depende de la complejidad del trabajo específico de la fase. La 
adaptación al Proyecto Integrador conserva las cuatro fases del 
modelo original y les asocia una etapa del proyecto, un conjunto de 
artefactos de cierre y mecanismos de retroalimentación entre fases. 
El esquema resultante se ilustra en la Fig.~\ref{fig:diagrama-rup}.

\begin{figure}[H]
\centering
\includegraphics[width=\columnwidth]{diagrama_rup.png}
\caption{Adaptación de las cuatro fases del RUP al ciclo de vida del Proyecto Integrador, con sus respectivas etapas, artefactos de cierre y retroalimentaciones entre fases.}
\label{fig:diagrama-rup}
\end{figure}

La fase de Inicio corresponde a la etapa de definición del problema y 
alcance del proyecto, y produce como artefactos de cierre el acta de 
constitución y el documento de definición. La fase de Elaboración 
corresponde a la etapa de análisis de requerimientos y planificación, 
y cierra con el documento de requerimientos y los diagramas de 
arquitectura. La fase de Construcción corresponde a la etapa de 
diseño y desarrollo del producto, y produce las evidencias de pruebas 
y el código fuente de los incrementos. La fase de Transición 
corresponde a la etapa de divulgación de resultados del proyecto y 
del producto, y produce el informe de pruebas y despliegue junto con 
el informe final.

Entre fases operan tres flujos de retroalimentación. Los ajustes de 
alcance originados en Elaboración pueden afectar los supuestos 
establecidos durante Inicio. Los riesgos y ajustes surgidos durante 
Construcción retroalimentan las decisiones de Elaboración. Los 
hallazgos del proceso de retroalimentación y los cambios identificados 
en etapas finales se reinyectan hacia fases tempranas, garantizando 
que los artefactos producidos al cierre de cada hito reflejen el 
estado real del desarrollo y no una especificación congelada al 
inicio del proyecto.

\subsection{Fases, hitos y artefactos}

La fase de \textit{Inicio} establece el caso de negocio y delimita el 
sistema mediante la identificación del problema, las causas que lo 
originan y los actores involucrados. Los artefactos producidos son el 
acta de constitución del proyecto, que fija el alcance y las 
exclusiones, y el documento de definición, que desarrolla el árbol 
del problema, los objetivos y la primera revisión del estado del arte 
que justifica la viabilidad técnica del enfoque propuesto.

La fase de \textit{Elaboración} traduce los objetivos en 
requerimientos verificables y establece la estructura técnica del 
sistema. Los requerimientos funcionales y no funcionales se 
documentan siguiendo el estándar IEEE 830, con identificadores 
únicos, prioridades y criterios de aceptación que habilitan la 
trazabilidad posterior hacia los casos de prueba. En paralelo se 
construyen los diagramas de arquitectura C4 (contexto, contenedores, 
componentes) que describen la descomposición del sistema y las 
interfaces entre componentes, junto con la matriz de riesgos que 
califica cada riesgo identificado en términos de probabilidad e 
impacto.

La fase de \textit{Construcción} es la más extensa del ciclo y se 
ejecuta de forma iterativa. Cada ciclo produce un incremento 
funcional del sistema: primero el núcleo del sensor de captura con 
su canal TCP y el receptor correspondiente; luego el canal UDP de 
métricas y el motor de persistencia en PostgreSQL; después el 
pipeline de análisis con Zeek y el motor de detección de anomalías; 
posteriormente la API REST, el canal WebSocket y la aplicación móvil; 
y finalmente la configuración de la infraestructura de integración y 
entrega continua. Cada iteración se cierra con pruebas unitarias e 
integración del incremento producido, que en conjunto conforman las 
evidencias de prueba y el código fuente versionado de los 
incrementos.

La fase de \textit{Transición} evalúa el éxito del proyecto mediante 
la comparación directa con los objetivos establecidos: revisión de 
los entregables, contrastación entre funcionalidades implementadas y 
requerimientos especificados, validación de las pruebas realizadas y 
medición del cumplimiento del cronograma. El resultado se consolida 
en el informe de pruebas y despliegue, y culmina con la producción 
del informe de investigación y su sustentación ante 
la Jornada de Socialización de Proyectos Integradores de la Facultad 
de Ingeniería de Sistemas e Informática de la Universidad Pontificia 
Bolivariana.

\subsection{Herramientas de gestión, control de versiones e integración continua}

Las fases descritas en las secciones anteriores requieren un sustrato
operativo que preserve la trazabilidad de cada decisión, mantenga el
estado compartido del código entre los integrantes del equipo y provea
retroalimentación automatizada sobre la calidad de cada incremento.
Este sustrato se articula mediante un sistema de control de versiones
distribuido, una plataforma de gestión del trabajo alineada con los
requerimientos funcionales, y un conjunto de pipelines de integración
y entrega continua que ejecutan sobre cada cambio las verificaciones
necesarias antes de aceptarlo como parte del producto en
construcción.

\subsubsection{Control de versiones y gestión del trabajo}

El sistema se mantiene en dos repositorios Git alojados en GitHub,
uno para el backend (\texttt{goatguard-server}) y otro para la
aplicación móvil (\texttt{goatguard-app}). La separación responde a
la divergencia de ciclos de vida, entornos de ejecución y
dependencias tecnológicas de ambos componentes: el backend evoluciona
en Python sobre contenedores Linux mientras que la aplicación móvil
lo hace en Dart sobre empaquetamiento Android. Mantenerlos en
repositorios distintos permite que cada uno progrese de forma
independiente y habilita pipelines de validación especializados para
cada stack. El código del agente de captura, por compartir el
runtime Python con el backend y por su naturaleza auxiliar respecto
al colector, se conserva dentro del repositorio del servidor para
preservar la coherencia del canal binario descrito en la
Sección~\ref{subsec:sensor}.

La estrategia de ramas adoptada es de tipo \textit{trunk-based}:
existe una única rama principal \texttt{main} que representa en todo
momento el estado aceptado del sistema, y los cambios se integran
mediante \textit{pull requests} sujetos a revisión cruzada entre los
tres integrantes del equipo antes de ser fusionados. Esta política,
acordada como práctica del equipo más que impuesta por reglas
automatizadas del repositorio, cumple la función de auditoría técnica
previa al merge y distribuye el conocimiento del código entre los
desarrolladores, evitando la concentración del criterio sobre zonas
específicas del sistema.

La planificación de las iteraciones de Construcción se coordina
mediante un tablero Jira, en el que los requerimientos funcionales
se descomponen en historias y tareas con trazabilidad hacia los
identificadores RF correspondientes. Jira complementa al histórico
de Git: este último registra la evolución del código, mientras que
aquel conserva el razonamiento de alto nivel y el avance por
iteración.

\subsubsection{Integración y entrega continua como práctica
metodológica}

Los pipelines de integración y entrega continua no constituyen un
aspecto meramente operativo del proyecto sino una pieza metodológica
central, en el sentido planteado por Humble y
Farley~\cite{humble2010cd}: cada cambio sobre la rama principal
debe atravesar un conjunto de verificaciones automáticas en una
ventana acotada de tiempo, de modo que el costo de detectar una
regresión se mantenga proporcional a la magnitud del cambio que la
introdujo y no se acumule hasta el final de la iteración. En el
contexto del RUP adoptado por el proyecto~\cite{kruchten2003rup},
esta práctica materializa el mecanismo de retroalimentación entre
Construcción y Elaboración: cuando una prueba automatizada contradice
un supuesto consignado en la fase de Elaboración, la evidencia queda
registrada en el mismo commit que introduce el cambio, y el
artefacto correspondiente puede revisarse antes del siguiente hito.

La Tabla~\ref{tab:gestion-rup} sintetiza la correspondencia entre
las cuatro fases del RUP adaptado, los artefactos de gestión que
produce cada una y las herramientas concretas empleadas para
materializarlos. La descripción técnica detallada de los pipelines,
las suites de pruebas y los artefactos de despliegue se desarrolla
en la Sección~\ref{subsec:cicd}.

\begin{table}[ht]
\centering
\caption{Correspondencia entre fases del RUP adaptado, artefactos de gestión y herramientas.}
\label{tab:gestion-rup}
\footnotesize
\begin{tabular}{@{}p{1.8cm}p{2.8cm}p{2.5cm}@{}}
\toprule
\textbf{Fase} & \textbf{Artefacto de gestión} & \textbf{Herramienta} \\
\midrule
Inicio & Acta de constitución, definición del problema & Documentos versionados en repositorio \\[3pt]
Elaboración & Backlog priorizado, diagramas de arquitectura & Jira, repositorios Git \\[3pt]
Construcción & Commits, pull requests, artefactos de CI & GitHub, GitHub Actions \\[3pt]
Transición & Imágenes OCI, APK, informe final & GHCR, artefactos de GitHub Actions \\
\bottomrule
\end{tabular}
\end{table}